\section{The \texorpdfstring{$I$}{I}-economy}
\label{bdc:sec:economy}

The formulae of the previous section share a curious economy: although the
process evolves in time, each moment can be written as a function of $I$ alone.
This is more than a change of variable. It says that the stochastic dynamics
move along a family of algebraic curves already fixed by the Riccati equation.

On the interval $(b,a)$, the right-hand side of
\cref{bdc:eq:riccatifact} is negative. Beginning at $I(0)=1$, the function
$I(t)$ therefore decreases strictly towards $b$, giving a bijection from
$[0,\infty)$ to $(b,1]$. It may be used as a clock, albeit one which runs
backwards: once a quantity is known as a function of $I$, its dependence on time
is determined by the speed
\[
\ddt{I}=\lambda(I-a)(I-b).
\]

The first two load moments suggest that the polynomial structure continues, and
indeed it does.

\begin{proposition}[Every load moment is a polynomial in $I$]
\label{bdc:prop:polynomial}
For each $p\ge1$, there is a polynomial $Q_p$ such that
\[
\ex{X_t^p}=Q_p\bigl(I(t)\bigr).
\]
It has degree $p+1$, satisfies $Q_p(a)=Q_p(b)=0$, and has leading coefficient
$(-1)^p p!(\lambda/\delta)^p$.
\end{proposition}

\begin{proof}
Write $m_p(t)=\ex{X_t^p}$. Applying the generator associated with
\cref{bdc:eq:rates} to $x\mapsto x^p$ gives
\[
\dot m_p
=-\delta m_{p+1}
+\bigl(\text{a linear combination of }m_1,\ldots,m_p\bigr).
\]
The highest moment comes from catastrophe: $X_t^p$ is replaced by zero at rate
$\delta X_t$. Solving for $m_{p+1}$ expresses it in terms of $\dot m_p$ and
lower moments. Along the Riccati trajectory,
\[
\frac{\dd}{\dd t}
=\lambda(I-a)(I-b)\frac{\dd}{\dd I},
\]
so, if $m_p=Q_p(I)$,
\[
m_{p+1}
=-\frac{\lambda}{\delta}(I-a)(I-b)Q_p'(I)
+\frac{1}{\delta}
\bigl(\text{a linear combination of }Q_1,\ldots,Q_p\bigr).
\]
By induction, every term contains the factor $(I-a)(I-b)$. The derivative term
raises the degree from $p+1$ to $p+2$ and multiplies the leading coefficient by
$-(\lambda/\delta)(p+1)$; the remaining terms have degree at most $p+1$.
The base case
\[
Q_1(I)=J(I)=-\frac{\lambda}{\delta}(I-a)(I-b)
\]
has degree two and leading coefficient $-\lambda/\delta$. The stated degree,
roots and coefficient follow.
\end{proof}

\FloatBarrier

Each new moment is therefore obtained by the same operation: differentiate the
previous polynomial along the Riccati flow, then add lower-order terms from
births and deaths. The cases $p=1$ and $p=2$ are precisely
\cref{bdc:eq:J,bdc:eq:KI}.

The variance retains this algebraic dependence. Substituting
\cref{bdc:eq:J,bdc:eq:KI} into \cref{bdc:eq:VarX} gives
\begin{equation}
\label{bdc:eq:VarXI}
\var{X_t}
=\frac{2\lambda^2}{\delta^2}(I-\kappa)(I-a)(I-b)
-\Bigl[\frac{\lambda}{\delta}(I-a)(I-b)\Bigr]^2.
\end{equation}
This quartic has the common factor $(I-a)(I-b)$; the two roots are generally
simple, not double. It also vanishes at $I=1$, as it must because $X_0=1$ is
deterministic. Along the part of the curve that has probabilistic meaning,
$I\in(b,1]$, the variance is positive for $t>0$ and tends back to zero as
$I\downarrow b$. Values on the algebraic continuation from $1$ towards $a$ are
not visited by the process and need not admit a probabilistic interpretation.

The release variance is likewise a polynomial in $I$: substitute
\cref{bdc:eq:V,bdc:eq:KI} into \cref{bdc:eq:VarW}. Its expanded form is long but
adds no new structure, so we leave it factored through the quantities already
defined.

\begin{figure}[!htbp]
\centering
\includegraphics[width=0.72\textwidth]{F4_I_economy.pdf}
\caption{The mean load $J$, mean release $V$ and second load moment $K$ as
functions of $I$, for $\lambda=1$, $\mu=0.2$ and $\delta=0.05$. The trajectory
runs from $I=1$ towards $b$; the symmetric-logarithmic ordinate keeps the three
scales visible without changing their values.}
\label{bdc:fig:economy}
\end{figure}

\Cref{bdc:fig:economy} removes time from the horizontal axis. The curves are
fixed algebraically, while the Riccati equation determines how quickly they are
traversed. This representation becomes useful in
\fwd{population}{Chapter~6}, where the single-cell formulae act as
age-dependent kernels in a renewal system. Since those kernels are polynomials
in one monotone coordinate, they can be manipulated analytically. If a
higher-order kernel is needed, \cref{bdc:prop:polynomial} supplies it by the same
recursion rather than by a fresh closure argument.

\FloatBarrier
